\subsection{The \texttt{open()} Function and File Object Construction}

The \texttt{open()} function is the ordinary entry point through which Python code obtains access to a file. It does not return the file contents directly. It returns a runtime object that represents an open communication channel between the Python process and an external filesystem resource.

That object is usually called a file object or stream object. It is a Python-level wrapper around lower-level operating-system access. The file object has methods such as \texttt{read()}, \texttt{write()}, \texttt{close()}, \texttt{flush()}, \texttt{seek()}, and \texttt{tell()}. Through these methods, Python code can move data across the boundary between runtime memory and persistent storage.

\subsubsection{Path Argument, Mode Argument, Encoding Argument, and Runtime File Object Creation}

The most common call shape is:

\begin{verbatim}
open(path, mode, encoding=...)
\end{verbatim}

The \texttt{path} argument identifies the external file. The \texttt{mode} argument tells Python what kind of access is requested. The \texttt{encoding} argument matters in text mode because stored bytes must be decoded into \texttt{str} objects when reading, or encoded from \texttt{str} objects when writing.

\begin{verbatim}
path = "io_demo_1.txt"
text = "It's just a flesh wound.\n"

file_obj = open(path, "w", encoding="utf-8")

print("The file object was created by open().")
print(f"path     = {path}")
print(f"text     = {text.strip()}")
print(f"file_obj = {file_obj}")

print()
print("Object types:")
print(f"type(path)     = {type(path)}")
print(f"type(text)     = {type(text)}")
print(f"type(file_obj) = {type(file_obj)}")

print()
print("Selected names from dir(file_obj):")
for name in ["write", "read", "close", "flush", "seek", "tell", "encoding", "mode"]:
    print(f"{name:10} -> {name in dir(file_obj)}")

print()
written_count = file_obj.write(text)
file_obj.close()

print(f"write() returned: {written_count}")
print("The returned number is the number of text characters written.")
print("The file now exists outside the Python runtime.")
\end{verbatim}

The name \texttt{file\_obj} does not contain the file itself. It refers to a Python object that manages access to the external file. The file object knows its mode, its encoding in text mode, and the operations allowed by that mode.

The call to \texttt{write()} receives a \texttt{str} object. Because the file was opened in text mode with \texttt{encoding="utf-8"}, Python encodes that string into bytes before they are stored externally. The call to \texttt{close()} releases the open file resource. Resource lifetime will be examined more carefully in the next subsection.

\subsubsection{Read Modes, Write Modes, Append Modes, and Exclusive Creation Modes}

The mode string controls the intended access pattern. The basic modes are:

\begin{itemize}
    \item \texttt{"r"}: read an existing file;
    \item \texttt{"w"}: write a file, replacing existing contents;
    \item \texttt{"a"}: append to the end of a file;
    \item \texttt{"x"}: create a new file, failing if it already exists.
\end{itemize}

These modes are not cosmetic. They change what the file object is allowed to do and what happens to existing stored data.

\begin{verbatim}
path = "io_modes_demo.txt"

file_obj = open(path, "w", encoding="utf-8")
file_obj.write("First line: D'oh!\n")
file_obj.close()

print("After write mode:")
file_obj = open(path, "r", encoding="utf-8")
content = file_obj.read()
file_obj.close()
print(content)

file_obj = open(path, "a", encoding="utf-8")
file_obj.write("Second line: Serenity now!\n")
file_obj.close()

print("After append mode:")
file_obj = open(path, "r", encoding="utf-8")
content = file_obj.read()
file_obj.close()
print(content)

print("Type inspection:")
print(f"type(content) = {type(content)}")

print()
print("Selected names from dir(content):")
for name in ["splitlines", "upper", "replace", "strip", "count"]:
    print(f"{name:12} -> {name in dir(content)}")

print()
print("Trying exclusive creation on an existing file:")
try:
    file_obj = open(path, "x", encoding="utf-8")
    file_obj.write("This should not replace the old file.\n")
    file_obj.close()
except FileExistsError as err:
    print(f"FileExistsError: {err}")
\end{verbatim}

The first \texttt{"w"} opening creates or replaces the file. The later \texttt{"a"} opening preserves the existing data and adds new text at the end. The later \texttt{"r"} openings read the stored data back into Python as a \texttt{str} object.

The \texttt{"x"} mode is useful when overwriting would be a bug. It asks Python to create the file only if the path does not already exist. If the file already exists, Python raises \texttt{FileExistsError} instead of silently destroying previous contents.

\subsubsection{Text Mode vs. Binary Mode: \texttt{str} Streams vs. \texttt{bytes} Streams}

A file can be opened in text mode or binary mode.

Text mode works with \texttt{str} objects. Python performs encoding and decoding at the boundary. This is appropriate for ordinary textual data: source code, logs, CSV files, JSON files, configuration files, and plain text documents.

Binary mode works with \texttt{bytes} objects. Python does not interpret the bytes as text. This is appropriate for images, archives, compiled files, serialized binary formats, and any file where raw byte values matter.

The mode string controls this distinction. Modes such as \texttt{"r"}, \texttt{"w"}, and \texttt{"a"} are text modes by default. Modes such as \texttt{"rb"} and \texttt{"wb"} are binary modes.

\begin{verbatim}
text_path = "io_text_demo.txt"
bin_path = "io_binary_demo.bin"

text_msg = "The answer is 42, but the question is unstable.\n"

file_obj = open(text_path, "w", encoding="utf-8")
file_obj.write(text_msg)
file_obj.close()

file_obj = open(text_path, "r", encoding="utf-8")
text_data = file_obj.read()
file_obj.close()

print("Text mode result:")
print(f"text_data       = {text_data.strip()}")
print(f"type(text_data) = {type(text_data)}")

print()
byte_data = text_msg.encode("utf-8")

print("Manual encoding creates bytes:")
print(f"byte_data       = {byte_data}")
print(f"type(byte_data) = {type(byte_data)}")

print()
print("Selected names from dir(byte_data):")
for name in ["decode", "hex", "split", "replace", "startswith"]:
    print(f"{name:12} -> {name in dir(byte_data)}")

file_obj = open(bin_path, "wb")
file_obj.write(byte_data)
file_obj.close()

file_obj = open(bin_path, "rb")
raw_data = file_obj.read()
file_obj.close()

print()
print("Binary mode result:")
print(f"raw_data       = {raw_data}")
print(f"type(raw_data) = {type(raw_data)}")

print()
decoded_again = raw_data.decode("utf-8")
print("Decoded back into text:")
print(f"decoded_again       = {decoded_again.strip()}")
print(f"type(decoded_again) = {type(decoded_again)}")
\end{verbatim}

In text mode, \texttt{read()} returns a \texttt{str}. In binary mode, \texttt{read()} returns a \texttt{bytes} object. This is the central difference.

The text file example lets Python handle the conversion between \texttt{str} and stored bytes through the declared encoding. The binary file example performs the conversion manually with \texttt{encode()} and \texttt{decode()}. This makes the boundary visible: external storage ultimately contains bytes, while Python text is represented as \texttt{str} objects inside the runtime.

The practical rule is direct: use text mode when the file represents text, and always specify an encoding such as \texttt{"utf-8"}. Use binary mode when the file represents raw bytes and must not be decoded as text.

